\exercisesfor{4}

在下列习题中，字母 G 表示一个群。

\begin{enumerate}
\def\labelenumi{\arabic{enumi}.}
\item
  设 \((\mathbf{G}_n)\) 是 \(G\) 上的一个整中心滤过。证明李代数 \(\mathrm{gr}(G)\) 由 \(\mathrm{gr}_1(G)\) 生成，当且仅当 \(\mathbf{G}_n = \mathbf{G}_{n+1}.\mathbf{C}^n\mathbf{G}\) 对所有 \(n \geqslant 1\) 成立。在此情形，证明 \(\mathbf{G}_n = \mathbf{G}_m.\mathbf{C}^n\mathbf{G}\) 对 \(m > n\) 成立，并推出 \((\mathbf{G}_n) = (\mathbf{C}^n\mathbf{G})\)，若存在整数 \(m\) 使 \(\mathbf{G}_m = \{\varepsilon\}\)。
\item
  设 \((\mathbf{G}_{\alpha})\) 是群 \(\mathbf{G}\) 上的实滤过，\(v\) 是相应的阶函数。令 \(\mathbf{H}\) 是 \(\mathbf{G}\) 的一个子群。
\end{enumerate}

\begin{enumerate}
\def\labelenumi{(\alph{enumi})}
\item
  令 \(H_{\alpha} = H \cap G_{\alpha}\)。证明 \((\mathrm{H}_{\alpha})\) 是 H 上的实滤过，且相应的阶函数是 v 在 H 上的限制。\(H_{\alpha}^{+} = H \cap G_{\alpha}^{+}\)，并且 \(\mathrm{gr}(\mathrm{H})\) 可与 \(\mathrm{gr}(\mathrm{G})\) 的一个分次子群相识别。
\item
  假设 H 是正规子群，且 \(v(G) \cap \mathbf{R}\) 是 R 的离散子集。记 \((G/H)_{\alpha} = (G_{\alpha}H)/H\)。证明 \(((G/H)_{\alpha})\) 是 G/H 上的实滤过，且相应的阶函数 \(v_{G/H}\) 由下式给出
\end{enumerate}

\[
v _ {\mathrm{G/H}} (x) = \operatorname * {S u p} _ {y \in x} v (y).
\]

证明对所有 \(\alpha \in \mathbf{R}\)，存在正合列

\[
0 \rightarrow \mathrm{gr} _ {\alpha} (\mathrm{H}) \rightarrow \mathrm{gr} _ {\alpha} (\mathrm{G}) \rightarrow \mathrm{gr} _ {\alpha} (\mathrm{G} / \mathrm{H}) \rightarrow 0.
\]

\begin{enumerate}
\def\labelenumi{(\alph{enumi})}
\setcounter{enumi}{2}
\tightlist
\item
  在 \((b)\) 的假设下，假设 \((G_{\alpha})\) 是中心滤过。证明由 \((G_{\alpha})\) 在 H 和 G/H 上诱导的滤过也都是中心滤过，并且李代数 \(\mathrm{gr}(\mathrm{G}/\mathrm{H})\) 可与 \(\mathrm{gr}(\mathrm{G})\) 对理想 \(\mathrm{gr}(\mathrm{H})\) 取商所得的代数相识别。
\end{enumerate}

¶ 3. 设 \((\mathrm{H}_{\alpha})\) 是群 H 上的实滤过，\(v_{\mathrm{H}}\) 是相应的阶函数。假设 G 在 H 上左作用，并且对所有 \(g \in G\)，映射 \(h \mapsto g(h)\) 都是滤过群 H 的自同构。

\begin{enumerate}
\def\labelenumi{(\alph{enumi})}
\tightlist
\item
  若 \(\alpha \in \mathbf{R}\)，用 \(G_{\alpha}\) 表示满足下列条件的 \(g \in G\) 的集合：
\end{enumerate}

\[
v _ {\mathrm{H}} (h ^ {- 1} g (h)) \geqslant v _ {\mathrm{H}} (h) + \alpha \quad \text { for   all } h \in \mathrm{H}.
\]

证明 \((\mathbf{G}_{\alpha})\) 是 \(\mathbf{G}\) 上的实滤过，并且 \(\mathbf{G}_{\alpha} = \mathbf{G}\) 当 \(\alpha \leqslant 0\) 时。令 \(v\) 表示相应的阶函数。

\begin{enumerate}
\def\labelenumi{(\alph{enumi})}
\setcounter{enumi}{1}
\tightlist
\item
  假设滤过 \((\mathrm{H}_{\alpha})\) 是中心滤过，且 \(\mathrm{G}_0^+ = \mathrm{G}\)。证明 \((\mathrm{G}_{\alpha})\) 是 G 上的中心滤过。若 \(\xi \in \mathrm{gr}_{\alpha}(\mathrm{G})\)、\(\eta \in \mathrm{gr}_{\beta}(\mathrm{H})\)，令 \(g\)（相应地，\(h\)）为 \(\xi\)（相应地，\(\eta\)）在 \(\mathrm{G}_{\alpha}\)（相应地，\(\mathrm{H}_{\beta}\)）中的代表元；证明 \(h^{-1}g(h)\) 在 \(\mathrm{gr}_{\alpha + \beta}(\mathrm{H})\) 中的像与 \(g\)、\(h\) 的选择无关；若将其记为 \(\mathrm{D}_{\xi}(\eta)\)，证明 \(\mathrm{D}_{\xi}\) 可延拓为 \(\mathrm{gr}(\mathrm{H})\) 的次数为 \(\alpha\) 的一个导子，并且 \(\xi \mapsto \mathrm{D}_{\xi}\) 是从李代数 \(\mathrm{gr}(\mathrm{G})\) 到 \(\mathrm{gr}(\mathrm{H})\) 的导子李代数的一个同态。若 \(v_{\mathrm{H}}(\mathrm{H}) \cap \mathbf{R}\) 包含于 \(\Gamma\)（\(\mathbf{R}\) 的离散子群），则 \(v(\mathrm{G}) \cap \mathbf{R}\) 也包含于该子群，并且对所有 \(\alpha \in \mathbf{R}\)，上面定义的映射 \(\xi \mapsto \mathrm{D}_{\xi}\) 都是单射。
\end{enumerate}

\begin{enumerate}
\def\labelenumi{\arabic{enumi}.}
\setcounter{enumi}{3}
\item
  设 H 是类为 c 的幂零群，G 是在 \(\mathrm{H}/(\mathrm{H},\mathrm{H})\) 上平凡作用的 H 的自同构群。证明 G 是类 \(\leqslant c - 1\) 的幂零群。（将习题 3 应用于 H 的下中心列，并注意 G 在 \(\mathrm{gr}(\mathrm{H})\) 上平凡作用。）证明若 H 是有限 p-群（p 为素数），则 G 也是如此（方法相同）。
\item
  设 \(\mathbf{K}\) 是一个交换域，\(\mathbf{L}\) 是 \(\mathbf{K}\) 的有限 Galois 扩张，Galois 群为 \(\mathbf{G}\)。令 \(v_{\mathbf{L}}\) 是 \(\mathbf{L}\) 的一个取值于 \(\mathbf{Z}\) 的赋值（交换代数，第六章，§3，第 2 号），且在 \(\mathbf{G}\) 下不变。若 \(g \in \mathbf{G}\)，记
\end{enumerate}

\[
v (g) = \sup _ {x \in L ^ {\bullet}} v _ {L} \left(\frac {g (x) - x}{x}\right).
\]

证明 v 是 G 上一个分离的整滤过 \((G_{n})\) 的阶函数，满足 \(G_{0} = G\)，且此滤过在 \(G_{1}\) 上的限制是中心滤过（应用习题 3，取 H 为 \(v_{L}\) 的理想）。证明 \(G_{1}\) 是 p-群（相应地，化为单位元），当 L 的剩余域的特征为 p \textgreater{} 0（相应地，为零）时；当域 L 和 K 具有相同的剩余域时，\(G_{1}\) 是交换代数第六章 §8 习题 11 (b) 中定义的同态 \(\phi\) 的核。

¶ 6. 令 K{[}G{]} 为 G 在 K 上的代数，I 为典范同态 K{[}G{]} → K 的核（“增广理想”）。于是


\(K[G] = K \oplus I\)，且 \(I\) 以 \(g - 1\) 为基，其中 \(g \in G - \{e\}\)。

\begin{enumerate}
\def\labelenumi{(\alph{enumi})}
\item
  若 \(n\) 是整数 \(\geqslant 0\)，令 \(I^n\) 表示 \(K[G]\) 的第 \(n\) 次幂 \(I\) 所成的理想。令 \(G_n\) 为由满足 \(g \in G\) 且 \(g - 1 \in I^n\) 的元素组成的集合。证明 \((G_n)\) 是 \(G\) 上的整中心滤过。特别地，\(G_n \supset C^n G\) 对所有 \(n\) 成立。
\item
  证明当 \(\mathbf{K} = \mathbf{Z}\) 时，映射 \(g \mapsto g - 1\) 在取商后定义了从 \(\mathrm{G} / (\mathrm{G}, \mathrm{G})\) 到 \(\mathrm{I} / \mathrm{I}^2\) 的同构。推出 \(\mathrm{G}_2 = \mathrm{C}^2\mathrm{G}.\dagger\)
\item
  设 K 是特征为零的域且 G 有限。证明 \(I^{n} = I\) 对所有 \(n \geqslant 1\) 成立。
\item
  设 K 是特征为 p \textgreater{} 0 的域且 G 是 p-群。证明当 n 充分大时 \(I^{n}\{0\}\)。（首先利用代数第一章 §6 第 5 号命题 11 证明每个单 K{[}G{]}-模都同构于 K；推出 I 是 K{[}G{]} 的根基，因而是幂零的，因为 K{[}G{]} 在 K 上有限秩。）
\item
  设 K = Z，且 G 具有如下性质：对每个 \(g \in G\)，若 \(g \neq 1\)，则存在素数 p、一个 p-群 P 以及同态 \(f: G \to P\)，使得 \(f(g) \neq e\)。证明在此情形下 \(\bigcap_{n} I^{n} = \{0\}\)。（立即将其归结为 G = P 的情形。将 (d) 应用于域 \(F_{p}\)，可知存在 m 使 \(I^{m} \subset p\)。\(Z[G]\)，而 I 是 Z{[}G{]} 中的直因子，所以这意味着 \(I^{m} \subset pI\)，从而 \(\bigcap_{n} I^{mn} \subset \bigcap_{n} p^{n}I\)；后者为 0，因为 I 是有限生成的 Abel 群。）
\end{enumerate}

\begin{enumerate}
\def\labelenumi{\arabic{enumi}.}
\setcounter{enumi}{6}
\item
  给 G 赋予滤过 \((\mathrm{C}^{\mathrm{n}}\mathrm{G})\)，并假设 \(\mathrm{gr}_{1}(\mathrm{G}) = \mathrm{G}/(\mathrm{G}, \mathrm{G})\) 是循环群。证明 \(\mathrm{gr}_{n}(\mathrm{G}) = \{0\}\) 对 \(n \geqslant 2\) 成立（使用命题 5），并推出 \(\mathrm{C}^{\mathrm{n}}\mathrm{G} = (\mathrm{G}, \mathrm{G})\) 对 \(n \geqslant 2\) 成立。
\item
  令 \(G = \mathbf{S}\mathbf{L}_{2}(\mathbf{Z})\)，并令 \(x = \begin{pmatrix} 1 & 1 \\ 0 & 1 \end{pmatrix}, y = \begin{pmatrix} 1 & 0 \\ 1 & 1 \end{pmatrix}, w = \begin{pmatrix} 0 & 1 \\ -1 & 0 \end{pmatrix}.\)
\end{enumerate}

\begin{enumerate}
\def\labelenumi{(\alph{enumi})}
\item
  验证公式 \(w^4 = 1\)、\(w = xy^{-1}x\)、\(wxw^{-1} = y^{-1}\)。
\item
  若 \(g = \begin{pmatrix} a & b \\ c & d \end{pmatrix}\) 是 \(G\) 的元素，记 \(l(g) = |a| + |b|\)。证明 \(l(g) = 1\) 当且仅当 \(g\) 具有 \(y^n w^a\) 的形式，其中 \(n \in \mathbf{Z}\) 且 \(0 \leqslant a \leqslant 3\)。若 \(l(g) \geqslant 2\)，证明存在某个幂 \(h\)（它是 \(x\) 或 \(y\) 的幂），使 \(l(gh) < l(g)\)。推出 \(G\) 由 \(\{x, y\}\) 生成。
\item
  利用 (a) 和 (b)，证明 G/(G, G) 由 x 的像 \(\xi\) 生成，且 \(\xi^{12} = e.\ddagger\) 推出 \(C^{n}G = (G, G)\) 对 \(n \geqslant 2\) 成立。（应用习题 7。）
\end{enumerate}

\begin{enumerate}
\def\labelenumi{\arabic{enumi}.}
\setcounter{enumi}{8}
\tightlist
\item
  给 G 赋予滤过 \((C^{n}G)\)。
\end{enumerate}

\begin{enumerate}
\def\labelenumi{(\alph{enumi})}
\item
  证明 \(\operatorname{gr}(\mathbf{G})\) 中由 \(\operatorname{gr}_2(\mathbf{G})\) 生成的理想是 \(\sum_{q\geqslant 2}\operatorname{gr}_q(\mathbf{G})\)。
\item
  令 \((x_{i})_{i\in I}\) 是 \(G\) 的生成族，并令 \(m\geqslant 1\)。假设对所有 \((i,j)\in \mathrm{I}^2\)，都有 \((x_i,x_j)^m\in \mathrm{C}^3\mathrm{G}\)。证明对所有 \(q\geqslant 2\) 及所有 \(u\in \mathrm{C}^q\mathrm{G}\)，都有 \(u^{m}\in \mathrm{C}^{q + 1}\mathrm{G}\)（使用 \((a)\)）。
\end{enumerate}

\begin{enumerate}
\def\labelenumi{\arabic{enumi}.}
\setcounter{enumi}{9}
\tightlist
\item
  设 \(x, y\) 属于 \(G\)，并令 \(r, s\) 是两个整数 \(\geqslant 1\)。假设 \((x^r, y^s) = e\)。（a）证明对所有 \(n \geqslant 1\)，都有 \((x, y)^{(rs)n} \in C^{n+2}G\)。（可以假设 \(G\) 由 \(\{x, y\}\) 生成；然后应用习题 9 (b)，注意 \((x^r, y^s) \equiv (x, y)^{rs} \mod C^3G\)。）
\end{enumerate}

\begin{enumerate}
\def\labelenumi{(\alph{enumi})}
\setcounter{enumi}{1}
\tightlist
\item
  假设 \(x^r = y^s = e\)；令 \(t\) 表示 \(r\) 和 \(s\) 的最大公约数。证明 \((x, y)^t \in \mathbf{C}^3\mathbf{G}\)，并推出 \((x, y)^{t^n} \in \mathbf{C}^{n+2}\mathbf{G}\) 对所有 \(n \geqslant 1\) 成立（方法相同）。
\end{enumerate}

\begin{enumerate}
\def\labelenumi{\arabic{enumi}.}
\setcounter{enumi}{10}
\tightlist
\item
  设 H 是 G 的一个子群，m 是满足 \(\geqslant 1\) 的整数。假设 G 由一个族 \((x_{i})_{i\in I}\) 生成，且对所有 i 都有 \(x_{i}^{m}\in H\)。
\end{enumerate}

\begin{enumerate}
\def\labelenumi{(\alph{enumi})}
\tightlist
\item
  令 H 带有由 G 上的滤过 (C \(^{n}\) G) 诱导的滤过，并令 gr(H) 与 gr(G) 的一个分次李子代数相识别，参见~习题 2。证明对所有 \(n \geqslant 0\)，
\end{enumerate}

\[
m ^ {n}. \mathrm{gr} _ {n} (\mathrm{G}) \subset \mathrm{gr} _ {n} (\mathrm{H}).
\]

推出对所有 \(z \in \mathbf{H}.\mathbf{C}^n\mathbf{G}\)，都有 \(z^m \in \mathbf{H}.\mathbf{C}^{n+1}\mathbf{G}\)。

\begin{enumerate}
\def\labelenumi{(\alph{enumi})}
\setcounter{enumi}{1}
\tightlist
\item
  若 G 是幂零群，证明存在整数 \(N \geqslant 0\)（只取决于 G 的幂零类），使得 \(z^{mN} \in H\) 对所有 \(z \in G\) 成立。
\end{enumerate}

\begin{enumerate}
\def\labelenumi{\arabic{enumi}.}
\setcounter{enumi}{11}
\item
  假设 G 是幂零群。令 H 是 G 的一个子群，m 是满足 \(\geqslant 1\) 的整数，且令 x、y 是 G 中满足 \(x^{m} \in H\) 和 \(y^{m} \in H\) 的元素。证明存在整数 \(N \geqslant 0\)，使得 \((xy)^{nN} \in H\)。（将习题 11 应用于由 \(\{x, y\}\) 生成的群及其与 H 的交。）
\item
  \begin{enumerate}
  \def\labelenumii{(\alph{enumii})}
  \tightlist
  \item
    令 \(\mathbf{F}\) 是具有两个元素的基本族 \(\{\overline{x},\overline{y}\}\) 的自由群，令 \(c\) 是整数 \(\geqslant 2\)，并令 \(m\) 是整数 \(\geqslant 1\)。记 \(\mathrm{F}^c = \mathrm{F} / C^c\mathrm{F}\)，并令 \(x,y\) 为 \(\overline{x},\overline{y}\) 在 \(\mathrm{F}^c\) 中的像。令 \(\mathrm{F}_m^c\) 为 \(\mathrm{F}^c\) 中由 \(\{x^{m},y\}\) 生成的子群。证明存在整数 \(\mathrm{N}\geqslant 0\)，使 \(z^{mN}\in \mathrm{F}_{m}^{c}\) 对所有 \(z\in \mathrm{F}^c\) 成立（使用习题 11）。
  \end{enumerate}
\end{enumerate}

\begin{enumerate}
\def\labelenumi{(\alph{enumi})}
\setcounter{enumi}{1}
\item
  令 \(\mathrm{I}^c\)（相应地，\(\mathbf{I}_m^c\)）为 \(\mathrm{F}^c\)（相应地，\(\mathbf{F}_m^c\)）的由 \(y\) 生成的正规子群。证明若 \(N\) 如上选取且 \(z \in \mathrm{I}^c\)，则 \(z^{n^k} \in \mathbf{I}_m^c\)（注意 \(\mathrm{F}^c / \mathrm{I}^c\) 是无限循环群，其生成元为 \(x\) 的像，并推出 \(\mathbf{F}_m^c / \mathbf{I}_m^c \to \mathbf{F}^c / \mathrm{I}^c\) 是单射，从而 \(\mathbf{I}_m^c = \mathbf{F}_m^c \cap \mathbf{I}^c\)）。特别地，\(xy^m x^{-1} \in \mathbf{I}_m^c\)。
\item
  假设 G 是幂零群。令 H 是 G 的一个子群，L 是 H 的一个正规子群。令 \(g \in G\)，并令 \(m \geqslant 1\) 使得 \(g^{m} \in H\)。证明当 N 充分大时，\(gl^{m^{N}} g^{-1} \in L\) 对所有 \(l \in L\) 成立。（若 G 的类 \textless c，取 (a) 中的 N，并使用同态 \(f: F^{c} \to G\)，满足 \(f(x) = g, f(y) = l\)；注意 \(f(F_{m}^{c}) \subset H, f(I_{m}^{c}) \subset L\)，然后应用上面的 (b)。）
\end{enumerate}

¶ 14. 设 P 是素数集合。若整数 n 满足 n ≠0 且其所有素因子都属于 P，则称 n 为 P-整数。若元素 \(x \in G\) 满足存在 P-整数 n 使 \(x^{n} = \epsilon\)，则称其为 P-挠元素；若 G 的每个元素都是 P-挠元素（相应地，除 e 外 G 中没有 P-挠元素），则称 G 为 P-挠群（相应地，P-无挠群）。若对所有 \(x \in G\) 和每个 P-整数 n，都存在 \(y \in G\) 使 \(x = y^{n}\)，则称 G 为 P-可除群。

假设 G 是幂零群。

\begin{enumerate}
\def\labelenumi{(\alph{enumi})}
\item
  令 H 是 G 的一个子群，令 \(H_{P}\) 为满足对 \(x \in G\) 存在 P-整数 n 使 \(x^{n} \in H\) 的元素所成的集合。证明 \(H_{P}\) 是 G 的子群（使用习题 12），且 \((\mathrm{H}_{\mathrm{P}})_{\mathrm{P}} = \mathrm{H}_{\mathrm{P}}\)。群 \(H_{P}\) 称为 H 在 G 中的 P-饱和。若 G 是 P-可除群，则 \(H_{P}\) 也是 P-可除群。
\item
  令 \(\mathbf{L}\) 是 \(\mathbf{H}\) 的正规子群。证明 \(\mathbf{L}_{\mathbb{P}}\) 是 \(\mathbf{H}_{\mathbb{P}}\) 的正规子群。（使用习题 13 (c) 证明：若 \(g \in \mathbf{H}_{\mathbb{P}}\) 且 \(l \in \mathbf{L}\)，则 \(glg^{-1} \in \mathbf{L}_{\mathbb{P}}\)。）
\end{enumerate}

特别地，若 L 是 G 的正规子群，则 \(L_{P}\) 也是正规子群，且 \(G/L_{P}\) 是 P-无挠的。推出 G 的 P-挠元素集合是使 G/N 为 P-无挠群的最小正规子群 N。

\begin{enumerate}
\def\labelenumi{(\alph{enumi})}
\setcounter{enumi}{2}
\item
  假设 G 是 P-无挠的。令 \(n\) 是 P-整数，且 \(x, y \in G\) 满足 \(x^n = y^n\)。证明 \(x = y\)。（对 \(r = s = n\) 应用习题 10 (a)。推出存在 P-整数 N 使 \((x, y)^{\mathbb{N}} = e\)，于是 \((x, y) = e\)，因为 G 是 P-无挠的；然后 \((x^{-1}y)^n = e\)，最终得到 \(x = y\)。）
\item
  令 H 是 G 的一个子群，L 是一个幂零群，且 \(f: H \to L\) 是同态。令 \(\Gamma\) 为 f 在 \(H \times L\) 中的图，令 \(\Gamma_{P}\) 为它在 \(G \times L\) 中的 P-饱和。证明 \(\Gamma_{P}\) 包含于 \(H_{P} \times L\)，并且当 L 是 P-可除群时 \(pr_{1}: \Gamma_{P} \to H_{P}\) 是满射，当 L 是 P-无挠群时是单射。
\end{enumerate}

假设 L 是 P-可除且 P-无挠的。证明 f 可唯一延拓为同态 \(f_{P}: H_{P} \to L\)，且 \(f_{P}\) 的图就是 \(\Gamma_{P}\)。

¶ 15. 沿用上题的记号，假设 G 是幂零群。令 \(i: \mathbf{G} \to \overline{\mathbf{G}}\) 是从 G 到幂零群 \(\overline{\mathbf{G}}\) 的同态。若满足下列条件，则称 \((i, \overline{\mathbf{G}})\) 为 G 的 P-包络：

\begin{enumerate}
\def\labelenumi{(\roman{enumi})}
\item
  \(\overline{\mathbf{G}}\) 是 P-可除且 P-无挠的。
\item
  \(i\) 的核是 G 的 P-挠元素集合。
\item
  \(i(\mathbf{G})\) 在 \(\overline{\mathbf{G}}\) 中的 P-饱和等于 \(\overline{\mathbf{G}}\)。
\end{enumerate}

\begin{enumerate}
\def\labelenumi{(\alph{enumi})}
\item
  令 \((i, \overline{\mathbf{G}})\) 是 G 的一个 P-包络，L 是 P-可除、P-无挠的幂零群。证明对每个同态 \(f\colon \mathbf{G} \to \mathbf{L}\)，存在唯一的同态 \(\bar{f}\colon \overline{\mathbf{G}} \to \mathbf{L}\)，满足 \(\bar{f} \circ i = f\)（将其归结为 G 是 P-无挠的情形，并使用习题 14 (d)）。推出若 G 有 P-包络，\(\dagger\) 则它在同构意义下唯一。
\item
  令 \((i, \overline{\mathbf{G}})\) 是 G 的一个 P-包络，令 H 是 G 的一个子群，令 \(\overline{\mathbf{H}}\) 是 \(i(\mathbf{H})\) 在 \(\overline{\mathbf{G}}\) 中的 P-包络，并令 \(i_{\mathbf{H}}: \mathbf{H} \to \overline{\mathbf{H}}\) 是由 \(i\) 诱导的同态。证明 \((i_{\mathbf{H}}, \overline{\mathbf{H}})\) 是 H 的一个 P-包络。
\end{enumerate}

假设 H 是 G 的正规子群。则 \(\bar{H}\) 是 \(\bar{G}\) 的正规子群，参见~习题 14 (b)；若 \(i_{G/H}: G/H \to \overline{G}/\overline{H}\) 是由 i 诱导的同态，证明 \((i_{G/H}, \overline{G}/\overline{H})\) 是 G/H 的一个 P-包络。

\begin{enumerate}
\def\labelenumi{\arabic{enumi}.}
\setcounter{enumi}{15}
\tightlist
\item
  沿用前两题的记号。
\end{enumerate}

\begin{enumerate}
\def\labelenumi{(\alph{enumi})}
\item
  令 \(\mathbf{S}_{\mathbf{P}}\) 为 P-整数的集合，令 \(\mathbf{\dot{Z}}_{\mathbf{P}} = \mathbf{S}_{\mathbf{P}}^{-1}\mathbf{Z}\) 为 \(\mathbf{Z}\) 关于 \(\mathbf{S}_{\mathbf{P}}\) 定义的分式环（交换代数，第二章，§2，第 1 号）。假设 G 是幂零、P-无挠且 P-可除的，并令 \(t \in \mathbf{Z}_{\mathbf{P}}\)、\(g \in G\)。证明存在唯一的属于 G 的元素 \(h\)，使得 \(h^s = g^{st}\) 对所有 \(s \in \mathbf{S}_{\mathbf{P}}\) 且 \(st \in \mathbf{Z}\) 的 s 成立。元素 \(h\) 称为 \(t\) 次 \(g\) 的幂，记为 \(g^t\)。映射 \(t \mapsto g^t\) 是从 \(\mathbf{Z}_{\mathbf{P}}\) 到 G 的同态。若 \(t\) 在 \(\mathbf{Z}_{\mathbf{P}}\) 中可逆，则 \(g \mapsto g^t\) 是双射。（使用习题 14 (c)。）
\item
  令 \(\mathbf{A}\) 是交换群，令 \(i\) 为从 \(\mathbf{A}\) 到 \(\mathbf{A}_{\mathbb{P}} = \mathbf{Z}_{\mathbb{P}} \times \mathbf{A}\) 的典范映射。证明 \((i, \mathbf{A}_{\mathbb{P}})\) 是 \(\mathbf{A}\) 的一个 P-包络。
\item
  令 \(\mathbf{G}\)（相应地，\(\overline{\mathbf{G}}\)）为 \(n\) 阶的环 \(k\)（相应地，环 \(k_{\mathrm{P}} = k\otimes \mathbf{Z}_{\mathrm{P}}\)）上的严格下三角群，并令 \(i\) 为从 \(\mathbf{G}\) 到 \(\overline{\mathbf{G}}\) 的由 \(k\) 到 \(k_{\mathrm{P}}\) 的典范同态所定义的同态（代数，第二章，§10，第 4 号）。证明 \((i,\overline{\mathbf{G}})\) 是 \(\mathbf{G}\) 的一个 P-包络。
\item
  令 G 是有限幂零群，因而是其 Sylow \(p\)-群 \(G_p\) 的直积（代数，第一章，§6，第 7 号，定理 4）。令 \(\overline{G}\) 为 \(G_p\)（\(p \notin P\)）的乘积，并令 \(i\) 为从 G 到 \(\overline{G}\) 的典范投影。证明 \((i, \overline{G})\) 是 G 的一个 P-包络。
\end{enumerate}

¶ 17. 假设 G 是幂零群。令 P 是素数集合，令 \(\mathfrak{V}_{P}(G)\) 为 G 中有限指数且指数为 P-整数的正规子群的集合（参见~习题 14）。令 \(\mathcal{T}_{P}(G)\) 为由滤子基 \(\mathfrak{V}_{P}(G)\) 在 G 上定义的拓扑（一般拓扑，第三章，§1，第 2 号，例）。

\begin{enumerate}
\def\labelenumi{(\alph{enumi})}
\item
  令 N 是 G 的一个有限指数子群，且 (G: N) 是 P-整数，令 \(N'\) 为 N 的所有共轭子群的交。证明 \((G: N')\) 是 P-整数（利用群 \(G/N'\) 是其 Sylow 子群的乘积这一事实）；推出 N 关于 \(\mathcal{T}_{\mathbb{P}}(G)\) 是开的。
\item
  令 H 是 G 的有限指数子群。证明 \(\mathcal{T}_{\mathrm{P}}(\mathrm{H})\) 与 \(\mathcal{T}_{\mathrm{P}}(\mathrm{G})\) 在 H 上诱导的拓扑重合。（方法相同。）
\item
  令 H 是 G 的正规子群，且 G/H 同构于 Z，令 x 是 G/H 的一个生成元在 G 中的代表元。令 \(N \in \mathfrak{V}_{\mathrm{P}}(\mathrm{H})\)，并令 \(N' = \bigcap_{n \in Z} x^{n} N x^{-1}\)。证明 \(N' \in \mathfrak{V}_{\mathrm{P}}(\mathrm{H})\)，且 \(x N' x^{-1} = N'\)。若 \(N''\) 是由 \(N'\) 和 x 生成的子群，证明 \((G: N'') = (H: N')\)。推出 \(\mathcal{T}_{\mathrm{P}}(\mathrm{H})\) 是由 \(\mathcal{T}_{\mathrm{P}}(\mathrm{G})\) 诱导的。
\item
  假设 \(\mathbf{G}\) 是有限群。证明若 \(\mathbf{H}\) 是 \(\mathbf{G}\) 的子群，则存在子群序列
\end{enumerate}

\[
\mathrm{H} = \mathrm{H} _ {0} \subset \mathrm{H} _ {1} \subset \dots \subset \mathrm{H} _ {k} = \mathrm{G}
\]

满足 \(\mathbf{H}_i\) 是 \(\mathrm{H}_{i + 1}\) 的正规子群，对 \(0\leqslant i < k\) 成立，且 \(\mathrm{H}_{i + 1} / \mathrm{H}_i\) 是有限群或

同构于 \(\mathbf{Z}\)。利用 \((b)\) 和 \((c)\)，推出 \(\mathcal{T}_{\mathrm{P}}(\mathrm{H})\) 是由 \(\mathcal{T}_{\mathrm{P}}(\mathrm{G})\) 诱导的。

\begin{enumerate}
\def\labelenumi{(\alph{enumi})}
\setcounter{enumi}{4}
\tightlist
\item
  在 (d) 的假设下，令 P' 表示素数集合中 P 的补集。证明 H 关于 \(\mathcal{T}_{\mathrm{P}}(\mathrm{G})\) 的闭包与 H 在 G 中的 P'-饱和重合（习题 14）；特别地，G 是 Hausdorff 的，当且仅当它是 P'-无挠的。
\end{enumerate}

\begin{enumerate}
\def\labelenumi{\arabic{enumi}.}
\setcounter{enumi}{17}
\tightlist
\item
  群 G 的上中心列 \((Z_{i}G)\) 递归定义如下：
\end{enumerate}

\begin{enumerate}
\def\labelenumi{(\roman{enumi})}
\item
  \(Z_{i}G = \{e\}\) 当 \(i \leqslant 0\) 时；
\item
  \(Z_{i}G/Z_{i-1}G\) 是 \(G/Z_{i-1}G\) 的中心。
\end{enumerate}

于是 \(\{e\} = Z_0G \subset Z_1G \subset \cdots\)，且 \(Z_1G\) 是 G 的中心；\(Z_iG\) 是 G 的特征子群。

\begin{enumerate}
\def\labelenumi{(\alph{enumi})}
\item
  证明 \(\mathbf{G}\) 是类 \(\leqslant c\) 的幂零群，当且仅当 \(\mathbf{G} = \mathbf{Z}_c\mathbf{G}\)。
\item
  证明 \((\mathbf{C}^n\mathbf{G},\bar{\mathbf{Z}}_m\mathbf{G})\subset \mathbf{Z}_{m - n}\mathbf{G}.\)
\item
  假设 G 是幂零且 P-无挠的（P 是素数集合，参见~习题 14）。证明 \(Z_{1}G\) 是 P-饱和的。（只需对 \(Z_{1}G\) 验证；若 n 是 P-整数，且 \(g \in G\) 满足 \(g^{n} \in Z_{1}G\)，则 \(xg^{n}x^{-1} = g^{n}\) 对所有 \(x \in G\) 成立，由习题 14 (c) 得 \(xgx^{-1} = g\)，从而 g 属于 \(Z_{1}G\)。）
\end{enumerate}

